\chapter{相关理论基础}\label{chap:related}\label{chap:theory}

本章只整理后文会直接用到的基础概念。多目标优化用于解释ECR与$J_{\min}$之间的取舍，粒子群优化对应第四章的MOPSO-DT求解器，计算几何则支撑复杂区域分解与坐标变换。这些理论共同服务于电子对抗场景下覆盖、压制与空地协同部署的多目标权衡。雷达探测、干扰强度和传播参数不在本章展开，统一放在第三章建模部分。

\section{多目标优化理论}

\subsection{从单目标到多目标优化}

在单目标优化中，候选方案通常可以按一个标量值排序；雷达部署并不符合这一简化。提高覆盖率可能需要把节点铺得更开，而提高最薄弱任务点的压制强度又可能要求节点向重点区域集中。两类目标方向不一致时，继续寻找“唯一最优解”意义不大，更合理的做法是保留一组互不支配的折中方案。

Pareto最优性最初由Vilfredo Pareto在经济学研究中提出，后来被进化多目标优化（Evolutionary Multi-objective Optimization, EMO）广泛采用。对本文而言，Pareto框架的作用很具体：它允许部署方案按照覆盖和压制两个目标同时比较，并把无法相互替代的方案保留下来，供后续任务偏好选择。

多目标优化通常同时关心两件事。第一是收敛性，即解集是否接近更优的目标区域；第二是多样性，即解是否覆盖前沿上的不同取舍位置。若粒子都集中在同一片目标空间，覆盖-压制曲线就会缺少可选方案；若只追求分散，又可能保留质量较差的解。因此，第四章的档案截断和领导者选择都会使用密度信息来约束解集分布。

\subsection{Pareto最优概念}

设候选解为$\mathbf{x}$，目标向量为$\mathbf{f}(\mathbf{x})=[f_1(\mathbf{x}),\ldots,f_m(\mathbf{x})]$。本文把各目标都写成最小化形式，因此目标值越小越优。

\begin{definition}[Pareto支配]
对于两个解$\mathbf{x}$和$\mathbf{y}$，若：
\begin{itemize}
    \item $\forall i \in \{1,2,\ldots,m\}, f_i(\mathbf{x}) \leq f_i(\mathbf{y})$（$\mathbf{x}$在所有目标上不差于$\mathbf{y}$）
    \item $\exists j \in \{1,2,\ldots,m\}, f_j(\mathbf{x}) < f_j(\mathbf{y})$（$\mathbf{x}$至少在一个目标上更优）
\end{itemize}
则称$\mathbf{x}$ Pareto支配$\mathbf{y}$，记为$\mathbf{x} \prec \mathbf{y}$。
\end{definition}

\begin{definition}[Pareto最优解]
若不存在任何可行解支配$\mathbf{x}^*$，则$\mathbf{x}^*$为Pareto最优解，也称非支配解。所有Pareto最优解映射到目标空间后形成Pareto前沿（Pareto Front, PF）。
\end{definition}

在本文实验中，一个非支配部署方案可能覆盖率较高但$J_{\min}$较低，另一个方案则相反。二者都可能进入档案，最终由任务偏好决定采用哪一个。

\subsection{多目标优化评价指标}

超体积（Hypervolume, HV）用于衡量非支配解集相对于参考点覆盖的目标空间体积。给定参考点$\mathbf{r} = (r_1, r_2)$，HV定义为：

\begin{equation}
\text{HV}(S) = \bigcup_{\mathbf{x} \in S} \text{vol}\left([f_1(\mathbf{x}), r_1] \times [f_2(\mathbf{x}), r_2]\right)
\label{eq:hv}
\end{equation}

HV同时反映收敛性和覆盖范围。若解集$A$严格支配解集$B$，则$\text{HV}(A) > \text{HV}(B)$，这一性质使它常被用作Pareto质量指标~\cite{Zitzler2001}。不过HV依赖参考点选择，本文在第五章统一参考点后再比较不同算法。

间距（Spacing）描述相邻非支配解之间距离是否均匀：

\begin{equation}
\text{Spacing}(S) = \sqrt{\frac{1}{|S|-1} \sum_{i=1}^{|S|} (\bar{d} - d_i)^2}
\label{eq:spacing}
\end{equation}

其中$d_i$为解$i$到最近邻的欧氏距离，$\bar{d}$为所有$d_i$的均值。Spacing越小，说明前沿上的候选方案分布越平滑；在部署决策中，这意味着任务人员能看到更连续的覆盖-压制取舍。

拥挤距离（Crowding Distance）估计解在前沿上的局部稀疏程度。NSGA-II中常用的计算方式为~\cite{Deb2002}：

\begin{equation}
d_i = \sum_{m=1}^{M} \frac{f_m(\mathbf{x}_{i+1}) - f_m(\mathbf{x}_{i-1})}{f_m^{\max} - f_m^{\min}}
\label{eq:crowding}
\end{equation}

拥挤距离较大的解位于稀疏区域，保留它们有助于补齐前沿。边界解通常被赋予$+\infty$，以免前沿两端的极端方案在截断时被删除。

\subsection{多样性维护机制的对比}

本文采用拥挤距离，但其他多目标算法也有自己的密度维护方式。这里列出三类在后续基线中会出现的方法。

\begin{itemize}
    \item SPEA2的$k$近邻密度估计根据第$k$近邻距离估计局部密度，常取$k=\sqrt{N+\bar{N}}$，其中$N$为种群规模，$\bar{N}$为归档规模。它不依赖坐标轴方向，对弯曲前沿较稳健~\cite{Zitzler2001}。
    \item 自适应网格把目标空间划分为网格，用每个网格中的解数量近似密度。该方法实现简单，但网格粒度会影响选择压力~\cite{Knowles2000}。
    \item MOEA/D的权重向量用于把多目标问题拆成多个标量子问题。权重向量分布越合理，得到的前沿覆盖通常越均匀~\cite{Zhang2007}。
\end{itemize}

拥挤距离的优点是无需额外密度参数，计算过程也容易嵌入MOPSO档案维护。由于本文只有ECR和$J_{\min}$两个物理目标，按目标轴排序得到的拥挤距离具有较直观的解释。

\subsection{标量化方法的简要讨论}

多目标问题也可以转成若干单目标问题求解，常见方法包括：

\begin{itemize}
    \item 加权和法：$\min \sum_i w_i f_i(\mathbf{x})$，$\sum_i w_i = 1$。实现最简单，但对非凸Pareto前沿覆盖不足。
    \item $\varepsilon$-约束法：保留一个目标，把其余目标改写为阈值约束。它能处理非凸前沿，但需要多次设定$\varepsilon$。
    \item 切比雪夫法：$\min \max_i \{w_i |f_i(\mathbf{x}) - z_i^*|\}$，其中$z_i^*$为参考点。该形式是MOEA/D中常用的子问题构造方式，理论上可逼近不同形状的前沿~\cite{Zhang2007}。
\end{itemize}

本文主算法仍采用Pareto支配和外部档案，因为一次运行即可得到多组部署候选解。MOEA/D则作为基线保留，用来观察分解式优化在同一部署模型下的表现。

\section{粒子群优化算法}

\subsection{标准粒子群优化}

粒子群优化（PSO）把每个候选解看作一个粒子，粒子位置表示当前解，速度控制下一步搜索方向，其基本形式最早提出于群智能优化研究中~\cite{Kennedy1995}。

\Needspace{8\baselineskip}
单目标PSO的连续变量更新为：
\vspace{-0.75\baselineskip}
\begin{align}
\mathbf{v}_i^{t+1}
&= w \cdot \mathbf{v}_i^t
 + c_1 \cdot r_1 \cdot (\mathbf{p}_i^{\text{best}} - \mathbf{x}_i^t) \nonumber\\
&\quad + c_2 \cdot r_2 \cdot (\mathbf{g}^{\text{best}} - \mathbf{x}_i^t) \label{eq:pso_velocity} \\
\mathbf{x}_i^{t+1}
&= \mathbf{x}_i^t + \mathbf{v}_i^{t+1} \label{eq:pso_position}
\end{align}
\vspace{-1.0\baselineskip}

其中$w$为惯性权重，$c_1$和$c_2$分别为认知因子和社会因子，$r_1,r_2\sim U(0,1)$为随机数，$\mathbf{p}_i^{\text{best}}$为粒子$i$的历史最优位置，$\mathbf{g}^{\text{best}}$为全局最优位置。第四章会在此基础上加入二进制区域编码和Pareto档案。

\subsection{惯性权重策略}

惯性权重$w$决定粒子延续上一轮速度的程度。$w$较大时，粒子更容易跨区域探索；$w$较小时，搜索更偏向局部收敛。本文比较三种设置：

\begin{enumerate}
    \item Legacy策略：$w(t) = 0.4 - 0.4 \cdot \frac{t}{T_{\max}}$，从0.4线性递减至0。
    \item Standard策略：$w(t) = 0.9 - 0.5 \cdot \frac{t}{T_{\max}}$，从0.9线性递减至0.4~\cite{Shi1998}。
    \item Adaptive策略：$w(t) = w_{\max} - (w_{\max} - w_{\min}) \cdot \left(\frac{t}{T_{\max}}\right)^{1/2}$，以非线性方式递减。
\end{enumerate}

第五章的参数实验显示，standard策略在本文场景中给出了较稳定的前沿分布，因此后续改进配置采用该策略。

\subsection{多目标PSO（MOPSO）}

把PSO扩展到多目标场景后，原来的“最优”概念需要重写。首先，当两个粒子互不支配时，个体最优不能只按单个目标判断；其次，档案中通常有多个非支配解，全局领导者不再唯一；最后，非支配解集规模会持续增长，需要截断机制。MOPSO通常用外部档案保存非支配解，再通过随机选择、网格选择或拥挤距离等方式挑选领导者。本文第四章采用拥挤度加权选择，以减少粒子在前沿密集区域反复聚集。

\section{计算几何基础}

\subsection{Hertel-Mehlhorn凸分解算法}

Hertel-\allowbreak Mehlhorn算法是一种简单多边形凸分解方法，时间复杂度为$O(n\log n)$，其中$n$为多边形顶点数~\cite{Hertel1983}。本文使用它的原因不是追求最少子区域数，而是希望把复杂区域拆成优化器容易编码和映射的凸子区域。其处理流程概括为：

\begin{enumerate}
    \item 连通性处理：将输入多边形拆分为独立连通分量。
    \item 空洞消除：对带空洞区域使用切割线把空洞转化为外边界的一部分。
    \item 三角剖分：对无空洞简单多边形进行Delaunay三角剖分。
    \item 三角形合并：在保持凸性的条件下贪心合并相邻三角形，得到凸多边形集合。
\end{enumerate}

分解后的每个凸子区域会获得一个二进制编码。粒子的离散变量负责选择子区域，连续变量负责决定该区域内部的部署坐标。

\subsection{坐标变换}

垂直交点坐标变换把归一化坐标$(\hx, \hy)\in[0,1]^2$映射到凸多边形$P$内的物理坐标$(x,y)$。本文采用的流程为：

\begin{enumerate}
    \item 根据$\hx$确定凸多边形包围盒中的垂直扫描线。
    \item 计算扫描线与多边形边界的上、下交点。
    \item 根据$\hy$在有效线段内插值得到物理纵坐标。
    \item 返回位于$P$内的物理坐标$(x,y)$。
\end{enumerate}

这种变换使粒子在归一化空间中的每次更新都能对应到合法区域内的位置，避免“更新后出界、再投影或删除”的处理。它不保证面积意义上的严格均匀采样；本文使用它主要是为了可行性和边界可达性，这一点在第四章算法说明和第五章边界实验中继续讨论。

\paragraph{评价指标的使用边界}
在本文的实验分析中，超体积、间距和 Pareto 前沿分布主要用于评价算法输出解集的整体质量，而不是直接替代雷达网络的物理性能指标。超体积可以同时反映解集向理想区域推进的程度和覆盖范围，但其数值会受到参考点选取、目标归一化方式以及样本规模的影响；间距指标可以刻画非支配解在目标空间中的均匀程度，但不能单独说明某个部署方案在工程上是否可用。因此，本文在解释实验结果时采用“算法评价指标”和“物理性能指标”并行的方式：前者用于比较不同优化策略的搜索表现，后者用于说明部署方案对应的探测代价和最小观测性能。这样的区分有助于避免把内部优化目标直接等同于最终工程指标，也使后续实验章中对 $J_{\min}$、ECR、HV 和 Spacing 的解释保持一致。

\paragraph{多目标评价与工程决策的关系}
多功能雷达网络部署问题通常不存在唯一的全局最优方案。若只追求较低的部署代价，算法可能倾向于选择数量较少或位置集中的雷达节点，从而降低边界区域或不利观测方向上的可观测性；若只追求较高的最小观测性能，则系统往往需要增加节点数量或扩大部署范围，带来更高的资源消耗。多目标优化的作用不是直接替决策者给出唯一答案，而是提供一组具有不同代价和性能权衡的候选方案。本文采用 Pareto 解集作为实验输出，一方面有助于保留不同部署偏好的可选空间，另一方面也便于观察区域分解、坐标变换和粒子群搜索策略对整体前沿形态的影响。

\section{本章小结}

本章梳理了本文后续建模和实验所需的基础概念。Pareto支配、HV、Spacing和拥挤距离用于描述非支配解集的质量；PSO和MOPSO提供连续变量搜索及外部档案维护框架；Hertel-Mehlhorn凸分解和垂直交点坐标变换则为复杂区域中的可行位置编码提供几何基础。下一章将在这些概念基础上给出面向电子对抗任务的空地协同部署问题形式化模型。
